\problemname{Sistakakan}
Katastrof! Joshua och Olle vill ha en kaka, men det finns bara en kvar i burken. Istället för att dela den på två föreslår Joshua att slå tärning om vem som ska få kakan.
Han tar fram $N$ stycken tärningar, vardera numrerade från $1$ till $M$. Tack vare att Joshua lärt sig lite geometri har han bevisat att alla sidor är lika sannolika att rullas.
Han låter dig nu välja $K$ stycken siffror, och om siffrorna på tärningarna summerar till ett av de talen vinner Olle, annars vinner Joshua.
Olle vill nu ha din hjälp att skriva ett program för att avgöra hans chans att vinna om han väljer sina utfall optimalt.


\section*{Indata}
Den första raden innehåller de tre heltalen $N, M, K$ ($1 \le N \le 20, 1 \le M, K \le 10000$), antalet tärningar, hur många sidor vardera tärning har och hur många utfall Olle får välja.

\section*{Utdata}
Skriv ut ett decimaltal -- den största sannolikheten olle kan ha att vinna om han väljer sina utfall rätt. 
Svaret kommer accepteras om de första 4 siffrorna efter kommatecknet stämmer överens med det korrekta svaret.

\section*{Poängsättning}
För ett poäng gäller att $n = 2, m \le 100$.\\
För ytterligare 1 poäng gäller att $n, m \le 6$.\\
För ytterligare 1 poäng gäller att $n, m \le 12$.\\
För ytterligare 1 poäng gäller att $n \le 20, n \le 100$.\\
För ytterligare 1 poäng gäller att $n \le 25, n \le 10000$.\\